\subsection{Binding, cost, and the Chapter 14 handoff}\label{sec:full:boundary}

\begin{definition}[Nonsuccinct full-relation soundness events]\label{def:full:sound-events}
For a cheating classical prover in the
\emph{full-relation mode}, introduce:
\begin{enumerate}[label=\textup{(\roman*)}]
\item $\mathcal B_{\rm pre}$:
the preprocessing proof accepts
but its knowledge extractor cannot
assign one source table $f$ and
mask $R$ to the commitments
\emph{before} the evaluation point;
\item $\mathcal B_{\rm rel}$:
an accepted relation proof
has a false circuit statement
for the fixed registered values;
\item $\mathcal B_{\rm bind}$:
one Pedersen bit commitment
has two distinct accepted
openings in the protocol;
\item $\mathcal B_{\rm hash}$:
one Merkle root has two
conflicting authenticated
commitment shells for the
same indexed position.
\end{enumerate}
Their probabilities are parameters
$\varepsilon_{\rm pre}$,
$\varepsilon_{\rm rel}$,
$\varepsilon_{\rm bind}$,
$\varepsilon_{\rm hash}$,
to be bounded by separate
cryptographic reductions in
Chapter~14.
\end{definition}

\begin{lemma}[Sound full-relation checks fix the original evaluation]\label{lem:full:fixed-eval}
Assume none of the events in
Definition~\ref{def:full:sound-events}
occurs and that the initial
claim relation
$\sigma_0=\sum_a f(a)\chi_0(a)=v$
is enforced using the \emph{same}
pre-evaluation extracted source
table $f$. Then
\[
 \boxed{v=\widetilde f(z).}
\]
\end{lemma}
\begin{proof}
The initial claim relation
identifies $v$ with
$\sum_a f(a)\chi_0(a)$.
By construction
$\chi_0(a)=\operatorname{eq}_m(a,z)$.
The multilinear interpolation
identity
Lemma~\ref{lem:intro:interpolation}
says exactly that this sum
equals $\widetilde f(z)$.
The no-bad-event hypotheses
ensure that the source table
in this relation is the same
one bound before $z$.
\end{proof}

\begin{theorem}[Conditional full-relation false-evaluation bound]\label{thm:full:sound-conditional}
For any adversary producing
an accepted claimed value
$v\neq\widetilde f(z)$
relative to the source table
extracted at preprocessing,
the acceptance probability
in full-relation mode is bounded by
\[
 \boxed{
 \varepsilon_{\rm full,snd}
 \le\min\{1,\,
 \varepsilon_{\rm pre}
 +\varepsilon_{\rm rel}
 +\varepsilon_{\rm bind}
 +\varepsilon_{\rm hash}\}.}
\]
This theorem does not
use a proximity reduction:
the entire word relation
is checked in the
nonsuccinct NP circuit.
\end{theorem}
\begin{proof}
By Lemma~\ref{lem:full:fixed-eval},
an accepted transcript outside
the four bad events must satisfy
$v=\widetilde f(z)$ for the
preprocessing-extracted table.
Thus acceptance of a false
evaluation implies at least
one of the four bad events.
The union bound gives their
total error, capped at one.
No sparse-query proximity
event has been invoked.
\end{proof}

\begin{definition}[Sparse KBFold proximity obligation]\label{def:full:proximity}
For the \emph{sparse} mode,
define $\varepsilon_{\rm prox}$
as a claimed bound on the event
that an adversary supplies
committed words outside one
consistent degree-bounded
masked fold chain but still
passes the authenticated
local checks, degree/proximity
tests, and hidden Sumcheck
relations. The event and bound
must be stated for the actual
sequence of enlarged Reed--Solomon
degree budgets
\[
 D_j=2^{m-j}+\lceil B/2^j\rceil
\]
and the correlation of its
adaptive fold and query challenges.
No numerical value is assigned
to $\varepsilon_{\rm prox}$
in this chapter.
\end{definition}

\begin{lemma}[Exact total fold-word size]\label{lem:full:total-words}
The total number of committed
field-symbol positions in
the entire $(m+1)$-level
word chain is
\[
 \boxed{
 \sum_{j=0}^{m}M_j
 =M(2-2^{-m})<2M.}
\]
\end{lemma}
\begin{proof}
By Lemma~\ref{lem:kbff:domain-sizes},
$M_j=M2^{-j}$. Summing the finite
geometric series gives
$M\sum_{j=0}^m2^{-j}
=M(2-2^{-m})$.
Because $m\ge1$, the last
factor is strictly less than two.
\end{proof}

\begin{lemma}[Bit-commitment count for fold oracles]\label{lem:full:word-count}
With the field encoding length
$b=\lceil\log_2|\F|\rceil$,
the full word chain uses
exactly
\[
 \boxed{bM(2-2^{-m})}
\]
Pedersen group elements for
its field-symbol commitments,
excluding separate source,
mask, round-polynomial and
offset commitments.
\end{lemma}
\begin{proof}
Every committed field symbol
has exactly $b$ bit commitments
by Definition~\ref{def:full:symbol-commit}.
Multiply the position count of
Lemma~\ref{lem:full:total-words}
by $b$.
\end{proof}

\begin{lemma}[Authenticated query-shell complexity]\label{lem:full:query-cost}
In a balanced binary Merkle tree
with $M_j$ leaves, an
authenticated commitment shell
contains $b$ group elements
and $\log_2 M_j$ sibling hashes,
apart from its public index.
Thus a local fold query
opens at most three such shells
and uses at most
\[
 \boxed{3b\text{ group elements}
       +3\log_2M\text{ sibling hashes}.}
\]
The zero-knowledge relation-proof
cost is \emph{additional}.
\end{lemma}
\begin{proof}
The tree height is
$\log_2M_j$ because $M_j$
is a power of two.
A path consists of one sibling
hash at each tree level.
The leaf holds $b$ group
elements. The fold query
requests two symbols at
level $j$ and one at $j+1$,
for three such shells.
Since $M_j\le M$, each
individual path has length
at most $\log_2M$, giving
the stated upper bound.
\end{proof}

\begin{proposition}[Full-relation mode has linear committed-word overhead]\label{prop:full:not-succinct}
The reference full-relation
mode requires $\Theta(bM)$
public group elements for
its authenticated word
commitments and at least
the full word arrays as
private relation witnesses.
It does not preserve the
polylogarithmic opening-size
objective of the source
KBFold PCS.
\end{proposition}
\begin{proof}
Lemma~\ref{lem:full:word-count}
gives a count between
$bM$ and $2bM$.
Consequently the public
commitment and input
representation alone are
linear in the expanded
word size $M$, up to
the encoding factor $b$.
A relation that checks the
entire word and its
Merkle commitments must
access each committed
leaf and its corresponding
hidden opening, so its
uncompressed witness
includes all these entries.
No compression theorem has
been supplied by this mode.
\end{proof}

\begin{definition}[Locked Chapter 14 security inputs]\label{def:full:handoff}
The following objects, and no
additional chapter, are the
Chapter~14 inputs:
\begin{enumerate}[label=\textup{(\roman*)}]
\item complete prover/verifier
$\Pi_{\rm KBFold}^{\rm hid}$,
its circuit relations and
online shell simulator;
\item the pre-evaluation
knowledge-binding reduction,
the Pedersen/hash consistency
errors and the online adaptive
relation-proof simulation contract;
\item the enlarged-budget sparse
proximity event
$\varepsilon_{\rm prox}$,
or the nonsuccinct
full-relation alternative;
\item the odd-characteristic
Sumcheck and linear fold-closure
identities;
\item a separate treatment of
classical interactive,
Fiat--Shamir ROM and
post-quantum/QROM models.
\end{enumerate}
Chapter~15 remains the
final Lean/Rust certification
and document audit block.
\end{definition}

\begin{remark}[Exact current status]
The algebraic prover/verifier
and its confidential
authentication checks are
fully defined and have
perfect completeness.
The Merkle/commitment
shell has exact witness-
independent online distribution.
A complete cryptographic
security theorem still needs
a concrete relation proof
engine with the precise
online simulation and
pre-evaluation extraction
guarantees, and the
sparse mode additionally
needs a proximity theorem
for the changed RS rates.
Those obligations are
confined to Chapter~14.
\end{remark}
